% !PRINTED 581
\weberpage{581}

%<<CONT
in which the $x_1, x_2, \ldots, x_n$ are integral rational numbers or integral rational functionals, one obtains an expression of the form
\weq{13}{\omega = A_0 + A_1\tau + A_2\tau^2 + \cdots + A_{n-1}\tau^{n-1},}
in which the $A_0, A_1, \ldots, A_{n-1}$ likewise denote \gesperrt{integral rational functionals}. By \S.\,145, 2., however, every integer and every integral functional of the field $\Omega$ is represented in this form, and we can therefore state the theorem:

3. \gesperrt{The powers}
\[
1, \tau, \tau^2, \ldots, \tau^{n-1}
\]
\gesperrt{form a basis of the integral functionals of the field $\Omega$.}

Instead of the powers of $\tau$, the functions
\[
F_0, F_1, F_2, \ldots, F_{n-1}
\]
can also be introduced as elements of the basis. These are defined by the equation
\weq{14}{\frac{F(t)}{t-\tau} = t^{n-1}F_0(\tau) + t^{n-2}F_1(\tau) + \cdots + F_{n-1}(\tau)}
and, as was already shown in \S.\,4 of the first Band, if
\[
F(t) = t^n + a_1t^{n-1} + \cdots + a_n
\]
they have the expression:
\weq{15}{\begin{aligned}
F_0(t) &= 1 \\
F_1(t) &= t + a_1 \\
F_2(t) &= t^2 + a_1 t + a_2 \\
&\cdots\cdots\cdots\cdots\cdots\cdots\cdots\cdots\cdots\cdots \\
F_{n-1}(t) &= t^{n-1} + a_1 t^{n-2} + a_2 t^{n-3} + \cdots + a_{n-1},
\end{aligned}}
so that the powers $1, t, t^2, \ldots, t^{n-1}$ can be represented without denominator by $F_0, F_1, \ldots, F_{n-1}$, and so that every integer or every integral functional $\omega$ also receives the expression:
\weq{16}{\omega = B_0 F_0(\tau) + B_1 F_1(\tau) + \cdots + B_{n-1} F_{n-1}(\tau),}
whose coefficients $B_0, B_1, \ldots, B_{n-1}$ are \gesperrt{integral rational functionals}.

If we now consider the integral function of $t$:
\[
S\,\frac{F(t)}{(t-\tau)\,F'(\tau)},
\]
% [sic: The German source prints a lowered comma after "so" at the line break; the German transcription retains it, while the English uses natural sentence punctuation.]
in which $S$ is the symbol for the trace explained in \S.\,134, it follows that it receives the value 1 for $t=\tau_1,\tau_2,\ldots,\tau_n$
%CONT>>
